\subsection{高斯投影测度的例子}

\begin{enumerate}
\def\labelenumi{\arabic{enumi})}
\tightlist
\item
  设 E 是实希尔伯特空间。映射 \(x' \mapsto \|x'\|^2\) 是 \(E'\) 上的正二次型。相应的高斯投影测度称为 E 上的典范高斯投影测度。可以证明，当 E 为无限维时，这个投影测度不是测度。
\end{enumerate}

设 \(A\) 是 \(\mathbf{E}\) 上的连续线性算子。映射 \(x' \mapsto \|{}^{t} A \cdot x'\|^2\) 是 \(\mathbf{E}'\) 上的正二次型。相应的投影测度 \(\mu_A\) 在 \(\mathbf{E}\) 上，当且仅当 \(A\) 是希尔伯特--施密特算子时才是测度（参见第~11号定理 3 的推论 2）。

\begin{enumerate}
\def\labelenumi{\arabic{enumi})}
\setcounter{enumi}{1}
\tightlist
\item
  正型核。设 T 是一个集合，\(E = R^{T}\) 是 T 上实函数构成的向量空间，并赋予逐点收敛拓扑。对于每个 \(t \in T\)，记 \(\varepsilon_{t}\) 为 E 上的线性形式 \(f \mapsto f(t)\)。族 \((\varepsilon_{t})_{t \in \mathbb{T}}\) 是 \(E'\) 的一个基（TVS, II, §6, 第~6号命题 8 的推论 2）。
\end{enumerate}

在 T 上，称每个满足下列关系的 \(T \times T\) 上的实值函数 K 为（实）正型核：

\[
\mathbf {K} (t, t ^ {\prime}) = \mathbf {K} (t ^ {\prime}, t) \quad \text { for } t, t ^ {\prime} \text { in } \mathbf {T},\tag{20}
\]

\[
\sum_ {i, j = 1} ^ {p} c _ {i} c _ {j} \mathrm{K} (t _ {i}, t _ {j}) \geqslant 0\tag{21}
\]

 对于任意正整数 \(p\)、元素 \(t_1, \ldots, t_p\)（属于 \(T\)）以及实数 \(c_1, \ldots, c_p\)。若上述条件成立，则公式

\[
q \left(\sum_ {t \in \mathrm{T}} c _ {t} \varepsilon_ {t}\right) = \sum_ {t, t ^ {\prime} \in \mathrm{T}} c _ {t} c _ {t ^ {\prime}} \mathrm{K} (t, t ^ {\prime})\tag{22}
\]

在 \(\mathbf{E}'\) 上定义一个正二次型。反过来，如果 \(q\) 是 \(\mathbf{E}'\) 上的正二次型，则公式

\[
\mathbf {K} (t, t ^ {\prime}) = \frac {1}{2} [ q (\varepsilon_ {t} + \varepsilon_ {t ^ {\prime}}) - q (\varepsilon_ {t}) - q (\varepsilon_ {t ^ {\prime}}) ]\tag{23}
\]

在 T 上定义一个正型核 K。由此，在 T 上的正型核集合与 \(E'\) 上的正二次型集合之间得到两个互为逆的双射。

设 K 是 T 上的正型核，q 是其在 \(E'\) 上的相关二次型。E 上方差为 q 的高斯投影测度也称为 E 上协方差为 K 的高斯投影测度。如果 T 可数，第~1号的命题 2 表明这个投影测度是测度。

\begin{enumerate}
\def\labelenumi{\arabic{enumi})}
\setcounter{enumi}{2}
\tightlist
\item
  设 \(\mathbf{T}\) 是可数集。通过令下式成立，定义正型核 \(\delta\)（在 \(\mathbf{T}\) 上）：
\end{enumerate}

\[
\delta (t, t ^ {\prime}) = \left\{ \begin{array}{l l} 1 & \text { if } t = t ^ {\prime} \\ 0 & \text { if } t \neq t ^ {\prime}. \end{array} \right.\tag{24}
\]

相应的二次型为 \(q\left(\sum_{t \in \mathbb{T}} c_t \varepsilon_t\right) = \sum_{t \in \mathbb{T}} c_t^2\)。对于每个 \(t \in \mathbb{T}\)，记 \(\mu_t\) 为 \(\mathbf{R}\) 上方差为 1 的高斯测度；

容易证明，\(\mathbf{R}^{\mathrm{T}}\) 上协方差为 \(\delta\) 的高斯测度等于 \(\bigotimes_{t\in T}\mu_t\)。

\begin{enumerate}
\def\labelenumi{\arabic{enumi})}
\setcounter{enumi}{3}
\tightlist
\item
  设 \(n \geqslant 1\) 是整数。若 \(C = (c_{ij})\) 是 \(n\) 阶方阵，且 \(\sum_{i,j=1}^{n} c_{ij} x_i x_j \geqslant 0\) 对任意实数 \(x_1, \ldots, x_n\) 都成立并且 C 对称，则称 C 为正对称矩阵；这等价于说映射 \((i,j) \mapsto c_{ij}\) 是集合 \(\{1,2,\ldots,n\}\) 上的正型核。因此可以谈论 \(\gamma_C\)，即 \(\mathbf{R}^n\) 上协方差为 \(C\) 的高斯测度；它由下式刻画
\end{enumerate}

\[
\int_ {\mathbf {R} ^ {n}} e ^ {i (x _ {1} t _ {1} + \dots + x _ {n} t _ {n})} d \gamma_ {C} (t _ {1}, \dots , t _ {n}) = \exp \left(- \frac {1}{2} \sum_ {j, k = 1} ^ {n} c _ {j k} x _ {j} x _ {k}\right),\tag{25}
\]

对于实数 \(x_{1}, \ldots, x_{n}\)。由第~5号的命题 6（公式 (19)），可推出

\[
\int_ {\mathbf {R} ^ {n}} t _ {j} t _ {k} d \gamma_ {C} (t _ {1}, \dots , t _ {n}) = c _ {j k} \quad (1 \leqslant j, k \leqslant n).\tag{26}
\]

由第~5号的命题 5，可推出公式

\[
u (\gamma_ {C}) = \gamma_{U C\,{}^{t}U}  ,\tag{27}
\]

其中 u 是从 \(R^{n}\) 到 \(R^{m}\) 的矩阵为 U 的线性映射。此外，很容易看出（参见第~5号命题 4 的证明的第一部分），如果 \(I_{n}\) 是 n 阶单位矩阵，则测度 \(\gamma_{I_{n}}\) 具有如下密度

\[
(2 \pi) ^ {- n / 2} \exp \left(- \frac {1}{2} (t _ {1} ^ {2} + \dots + t _ {n} ^ {2})\right)
\]

 相对于勒贝格测度 \(\lambda_{n}\)（在 \(\mathbf{R}^{n}\) 上）。

下面证明：如果矩阵 \(C\) 可逆，逆矩阵为 \(D = (d_{jk})\)，则

\[
\begin{array}{l} d \gamma_ {C} (t _ {1}, \dots , t _ {n}) = \\ (2 \pi) ^ {- n / 2} (\det D) ^ {1 / 2} \left(\exp \left(- \frac {1}{2} \sum_ {j, k = 1} ^ {n} d _ {j k} t _ {j} t _ {k}\right)\right) d t _ {1} \dots d t _ {n}. \end{array}\tag{28}
\]

 因为，如果 \(C\) 可逆，由下式定义的二次型 \(q\)（在 \(\mathbf{R}^n\) 上）

\[
q (x _ {1}, \dots , x _ {n}) = \sum_ {j, k = 1} ^ {n} c _ {j k} x _ {j} x _ {k}
\]

是非退化的。利用存在一个关于 q 正交规范的 \(R^{n}\) 基，可以证明存在 n 阶方阵 U，使得 \(C = U \cdot {}^{t}U\)，因而由 (27) 有 \(\gamma_{C} = u(\gamma_{I_{n}})\)（其中 u 表示矩阵为 U 的 \(R^{n}\) 自同构）。令 Q 为 \(R^{n}\) 上由下式定义的二次型

\[
Q (t _ {1}, \dots , t _ {n}) = t _ {1} ^ {2} + \dots + t _ {n} ^ {2};
\]

则

\[
\gamma_ {I _ {n}} = (2 \pi) ^ {- n / 2} e ^ {- Q / 2} \cdot \lambda_ {n},
\]

从而

\[
u (\gamma_ {I _ {n}}) = (2 \pi) ^ {- n / 2} e ^ {- (\mathrm{Q} \circ u ^ {- 1}) / 2} \cdot u (\lambda_ {n}).\tag{29}
\]

立即可见，\(\mathbf{Q} \circ u^{-1}\) 这个 \(\mathbf{R}^n\) 上的二次型的值为 \(\sum_{j,k=1}^{n} d_{jk} t_j t_k\)（在点 \((t_1, \ldots, t_n)\) 处），而第 VII 章第 1 节第~10号的命题 15 表明

\[
u (\lambda_ {n}) = (\det U) ^ {- 1} \cdot \lambda_ {n} = (\det D) ^ {1 / 2} \cdot \lambda_ {n}.\tag{30}
\]

公式 (28) 由此得到。

